\documentclass[11pt]{article}
\usepackage[a4paper,margin=1in]{geometry}
\usepackage{amsmath,amssymb,mathtools,bm}
\usepackage{enumitem,booktabs,array}
\usepackage{tcolorbox}
\usepackage{hyperref}
\usepackage[protrusion=true,expansion=false]{microtype}
\hypersetup{colorlinks=true,linkcolor=blue,urlcolor=blue,citecolor=blue}
\setlist[itemize]{topsep=2pt,itemsep=2pt}
\setlist[enumerate]{topsep=2pt,itemsep=2pt}
\newtcolorbox{keybox}{colback=blue!3!white,colframe=blue!40!black,boxrule=0.4pt,arc=1mm}
\newtcolorbox{alertbox}{colback=red!3!white,colframe=red!40!black,boxrule=0.4pt,arc=1mm}
\newtcolorbox{answerbox}{colback=green!3!white,colframe=green!40!black,boxrule=0.4pt,arc=1mm}
\newtcolorbox{drillbox}{colback=yellow!5!white,colframe=orange!60!black,boxrule=0.4pt,arc=1mm}
\newcommand{\EV}{\mathbb{E}}
\newcommand{\Var}{\mathrm{Var}}
\newcommand{\Cov}{\mathrm{Cov}}
\newcommand{\blank}[1]{\underline{\hspace{#1}}}

\title{E12-04: Michaely \& Moin (2022, JFE) \\
\large ``Disappearing and Reappearing Dividends'' --- Active Recall Workbook}
\author{Empirical Corporate Finance II --- Exam Prep}
\date{\today}
\begin{document}\maketitle

\begin{keybox}
\textbf{Paper in one sentence.} MM document that the long-run decline in US dividend-paying firms (Fama-French 2001's ``Disappearing Dividends'') stabilized and partially reversed post-2000: dividends reappeared not because firms returned to smoothing cash flows but because the universe of listed firms matured and repurchases + dividends converged to a unified ``total payout'' framework in which dividends play a smaller but meaningful share.

\textbf{Placement.} Updates Fama-French 2001 for the 21st century; important for payout-policy research in the era of rising buybacks.
\end{keybox}

\section*{Round 1: Complete Walk-Through}

\subsection*{1.1 Context}
Fama-French 2001 showed the fraction of public firms paying dividends fell from $\sim$70\% in 1978 to $\sim$20\% by 1999, driven by a surge of young unprofitable listings and a shift to repurchases. Did dividends continue to disappear?

\subsection*{1.2 Data}
\begin{itemize}
\item Compustat/CRSP, 1975--2019.
\item Dividend-paying status, payout levels, repurchases.
\item Firm characteristics: profitability, size, age, investment opportunities.
\end{itemize}

\subsection*{1.3 Facts}
\begin{itemize}
\item Dividend-paying fraction rose from $\sim$20\% (1999) to $\sim$35\% (2019).
\item Rise driven by older, profitable firms initiating dividends and sustained dividend continuations.
\item Aggregate dividend payout dollars grew steadily; repurchase dollars grew even faster.
\item Total payout (div + repurchase) / earnings trended up.
\end{itemize}

\subsection*{1.4 Analysis}
\begin{itemize}
\item Decompose the reappearance into composition effect (fewer young loss-making firms in sample) and propensity effect (existing firms more likely to initiate).
\item Roughly: 60\% composition, 40\% propensity.
\item Firms increasingly pair dividends with repurchases; total-payout framework best describes behavior.
\end{itemize}

\subsection*{1.5 Mechanism}
\begin{itemize}
\item Post-dot-com bust: fewer loss-making tech IPOs, sample composition shifts.
\item Mature firms in concentrated industries generate excess cash they prefer to return.
\item Dividend-tax cut (2003) may have mattered but composition shifts explain most of the reversal.
\end{itemize}

\subsection*{1.6 Implications}
\begin{itemize}
\item Dividend-policy research should account for sample composition.
\item Total-payout models fit better than dividend-only models.
\item Corporate-life-cycle theories (DeAngelo et al.\ 2006) gain support: mature firms pay out, young ones retain.
\end{itemize}

\subsection*{1.7 Caveats}
\begin{alertbox}
\begin{itemize}
\item Composition vs.\ propensity decomposition is sensitive to controls.
\item Private-to-public listing trends affect sample.
\item 2003 tax cut and other regulatory changes confound.
\end{itemize}
\end{alertbox}

\section*{Round 2: Level 1 Fill-in}
\begin{enumerate}
\item FF 2001 finding: \blank{2cm} dividends.
\item Dividend-paying fraction: 1999 $\sim$20\%, 2019 $\sim$\blank{1cm}\%.
\item Decomposition: $\sim$\blank{1cm}\% composition, $\sim$\blank{1cm}\% propensity.
\item Framework: \blank{2cm}-payout.
\item Life-cycle theory: DeAngelo \blank{1cm} 2006.
\end{enumerate}

\section*{Round 3: Level 2 Prompts}
\begin{enumerate}
\item Restate Fama-French 2001 and the reversal documented by MM.
\item Describe the composition-propensity decomposition.
\item Explain total-payout framework.
\item Discuss life-cycle implications.
\item Evaluate alternative explanations (tax cut, governance).
\end{enumerate}

\section*{Round 4: Minimal Scaffolding}
From a blank page: (i) context, (ii) facts, (iii) decomposition, (iv) framework, (v) implications.

\section*{Round 5: Blank Page}
Essay: dividends, repurchases, and the payout revival.

\vspace{6pt}
\begin{drillbox}
\textbf{Speed Drill.} (1) FF 2001 finding. (2) 2019 paying share. (3) Decomposition split. (4) Framework. (5) Life-cycle theory.
\end{drillbox}

\section*{Quick Reference \& Answer Key}
\begin{answerbox}
\begin{itemize}
\item \textbf{Context.} FF 2001 Disappearing Dividends.
\item \textbf{Fact.} Paying fraction 20\%$\to$35\%.
\item \textbf{Decomp.} 60\% composition, 40\% propensity.
\item \textbf{Model.} Total payout framework.
\end{itemize}
\end{answerbox}

\begin{keybox}
\textbf{30-second exam pitch.} Michaely \& Moin (2022) update Fama-French's (2001) ``Disappearing Dividends'' finding into the 21st century. Using Compustat/CRSP 1975--2019, the fraction of dividend-paying listed firms rose from $\sim$20\% in 1999 to $\sim$35\% in 2019, with aggregate dividend dollars growing steadily and repurchase dollars growing even faster. A decomposition attributes roughly 60\% of the recovery to composition shifts (fewer young loss-making listings post-dot-com) and 40\% to increased propensity of existing firms to initiate. Total payout / earnings trended up throughout. The paper endorses a total-payout framework (dividends + repurchases jointly) and supports corporate-life-cycle theories (DeAngelo et al.\ 2006): mature, profitable firms return cash; young ones retain. Dividend-research specifications should control for sample composition and use total-payout measures.
\end{keybox}

\end{document}
